\documentstyle[%


righttag%


,11pt%


,amssymb%


,russian%               remove in Eng.


,izvru%                 Set izven% in Eng.


% ,izvru-ks             for brief communications


]{amsart}





%


\newcommand{\tts}[1]{}





%\hoffset=-15mm%                  For Eng.


\setcounter{page}{71}


\god{2002}


\nomer{1~(476)} %                        Remove in Eng.


%\nomer{Vol.\,46, No.\,1}%               Set in Eng.


%\str{\pageref{begin}--\pageref{end}}%  Set in Eng.


\udk{515.124.4:514.774}





\author{\it Е.Н.\,Сосов}


% NB:   ^^ remove in Eng.


\title{О непрерывности метрической $\delta$-проекции на 


выпуклое множество в специальном метрическом


пространстве}


\thanks{Работа выполнена при финансовой поддержке Российского фонда


фундаментальных исследований, грант \No\,00-01-00308.}





\date{}


%\pagestyle{myheadings}%         Set in Eng.


\pagestyle{plain} %              Remove in Eng.


\begin{document}


%\thispagestyle{plain}%          Set in Eng.


\maketitle


\label{begin}








В статье доказано, что некоторые результаты 


А.В.\,Маринова из [1], [2] о непрерывности метрической


$\delta$-проекции на выпуклое множество в линейном 


нормированном пространстве остаются верными в


специальном метрическом пространстве.





\section{Необходимые определения и теоремы}





Пусть $(X,\rho )$ --- метрическое пространство


и $S$ --- выделенное в нем семейство сегментов


(т.\,е. кривых, длины которых равны расстояниями между их


концами [3], c.\,42), в дальнейшем


называемых хордами ([4], с. 23). 


Пространство $X$ и семейство хорд $S$ подчиним следующим условиям.


\begin{itemize}


\item[A.] Для каждой хорды все ее подсегменты (включая точки) 


являются хордами.


\item[B.] Каждые две точки пространства $X$ можно соединить 


единственной хордой.


\item[C.] Для всех $p, x,  y \in X$ выполняется неравенство 


$2\rho (p,\omega_{1/2}(x,y)) \leq \rho (p,x) + \rho (p,y)$,


где $\omega_{1/2}(x,y)$ --- середина хорды $[x,y]$ с концами $x$, $y$.


\end{itemize}





Известно ([4], с. 63; [3], с.\,304), что 


в прямых хордовых пространствах неположительной кривизны эти 


условия выполняются. Аналогичные локальные условия налагались 


на метрическое пространство в [5].


 


В дальнейшем используем следующие обозначения и определения:


$B[x,r]$ ($S(x,r)$) --- замкнутый 


шар (сфера) с центром в точке $x$, радиуса $r>0$;


$xy = \rho (x,y)$, $xM = \rho (x,M)$;


пусть $\delta \geq 0$, $x \in X$, $M \subset X$, $M \neq \emptyset$;


$x_{M}^{\delta} = \{y \in M : xy \le xM + \delta \}$


($ x_{M} = x_{M}^{0}$) --- метрическая


$\delta$-проекция (метрическая проекция) точки $x$ на множество $M$ [1];


$\beta (A,B) = \sup \{xB : x \in A \}$,


$\alpha (A,B) = \max \{\beta (A,B),\beta (B,A)\}$ ---


полууклонение и расстояние по Хаусдорфу между непустыми


множествами $A, B \subset X$ (бесконечные 


значения не исключаются) [1];


$\beta (x_{M}^{\delta \pm 0},F) =


\lim \limits_{t \to +0}{\beta (x_{M}^{\delta \pm t},F)}$,


где $N \subset X$ [1].


Множество $M$ пространства $X$, удовлетворяющего 


условиям A, B, называется выпуклым,


если для каждых двух точек $x$, $y$ из этого 


множества хорда $[x,y]$ принадлежит этому множеству.





Приведем два элементарных следствия условий A, B, C.





1. Для каждого $\lambda \in [0,1]$ и для всех $p$, $x$, $y$ 


из пространства $X$, удовлетворяющего условиям


A, B, C, выполняется неравенство


\begin{equation}


p\omega_{\lambda}(x,y) \leq (1 - \lambda )px + \lambda py,


\end{equation}


где точка $\omega_{\lambda}(x,y) \in [x,y]$ 


такая, что $x\omega_{\lambda}(x,y) = \lambda xy$.





2. Каждый замкнутый (открытый) шар пространства $X$, удовлетворяющего 


условиям A, B, C, является выпуклым множеством. 





Известна





\begin{lem}[{\series{m}\shape{n}\selectfont [1]}]


Пусть $X$ --- метрическое пространство, $x, y \in X$,


$F$, $M$, $N$ --- непустые множества из $X$, $\delta \geq 0$.


Тогда


$\beta (y_{N}^{\delta},F) \leq


\beta (N,M) + \beta (x_{M}^{\delta_{1} + 0},F)$,


$\beta (F,y_{N}^{\delta}) \leq \beta (M,N) +


\beta (F,x_{M}^{\delta_{2} - 0})$,


где


$\delta_{1} = \delta + xy + yN + \beta (N,M) - xM$,


$\delta_{2} =\delta + yN - xy - xM - \beta (M,N)$. Причем


предполагается, что $\delta_{2} > 0$.


\end{lem}





Сформулируем теперь полученные результаты. Oбобщает леммы 6--8 из [2] 





\begin{lem}


Пусть $M$ --- выпуклое множество пространства $X$, 


удовлетворяющего условиям {\em A, B, C;} $x \in X$,


$0 < t < \varepsilon < \delta $,


$0 < \varepsilon ' < \delta '$,


$\varepsilon ' \leq \varepsilon $, $\delta ' \leq \delta$


$(t \geq 0$ при $x_{M} \neq \emptyset)$,


$F \subset x_{M}^{\delta }$, $G \subset x_{M}^{t}$.


Тогда имеют место следующие неравенства$:$





{\em a)}  $\frac{\beta (F,x_{M}^{\varepsilon })}{\delta - \varepsilon }\leq


\frac{\beta (F,x_{M}^{t})}{\delta - t};$


{\em b)} $\frac{\beta (x_{M}^{\delta },G)}{\delta - t} \leq


\frac{\beta (x_{M}^{\varepsilon },G)}{\varepsilon - t};$


{\em c)} $\frac{\beta (x_{M}^{\delta },x_{M}^{\varepsilon })}{\delta -


\varepsilon } \leq \frac{\beta (x_{M}^{\delta' },


x_{M}^{\varepsilon' })}{\delta' - \varepsilon' }$.


\end{lem}





{\bf Следствие.}


Пусть $M$ --- выпуклое множество пространства $X$, 


удовлетворяющего условиям A, B, C; $x \in X$. Тогда


имеют место следующие неравенства:


\begin{itemize}


\item[a)] $\frac{\beta (x_{M}^{\delta },x_{M}^{+ 0})}{\delta } \leq


\frac{\beta (x_{M}^{\delta' },x_{M}^{+ 0})}{\delta' }$, где 


$\delta > \delta' > 0$;


\item[b)] $\beta (x_{M}^{\delta },x_{M}^{\varepsilon }) \leq


\frac{\beta (x_{M}^{\delta },x_{M}^{+ 0})}{\delta }(\delta - \varepsilon)$, 


где $\delta \geq \varepsilon > 0$;


\item[c)] $\alpha (x_{M}^{\delta },x_{M}^{\varepsilon }) \leq


\frac{\beta (x_{M}^{\mu },x_{M}^{+ 0}) }{\mu }|\delta - \varepsilon|


\leq \frac{\beta (x_{M}^{\delta },x_{M}^{+ 0}) }


{\delta }|\delta - \varepsilon|$, где $\delta > 0$, 


$\mu = \max (\delta,\varepsilon )$;


\item[d)] $\beta (x_{M}^{\delta },x_{M}^{\varepsilon }) \leq


(\frac{2xM}{\delta } + 1)(\delta - \varepsilon )$, 


где $\delta \geq \varepsilon \geq 0$.


При $x_{M} \neq \emptyset $ выражение $x_{M}^{+ 0}$ в этих


неравенствах можно заменить на $x_{M}$ (ср. [2], следствия 1--3).


\end{itemize}





Следующая лемма является обобщением леммы 5 из [2] 


и имеет аналогичное доказательство, если использовать


п.\,a) леммы 2.





\begin{lem}


Пусть $M$ --- выпуклое множество пространства $X$, 


удовлетворяющего условиям {\em A, B, C;} $x \in X$,


$\delta >0 $, $F \subset X$. Тогда имеют место следующие четыре


равенства$:$


$\beta (x_{M}^{\delta \pm 0},F) = \beta (x_{M}^{\delta },F)$,


$\beta (F,x_{M}^{\delta \pm 0}) = \beta (F,x_{M}^{\delta })$. 


\end{lem}





Aналогична лемме 10 из [2]





\begin{lem}


Пусть $M$ --- выпуклое множество пространства $X$, 


удовлетворяющего условиям {\em A, B, C;} $x, y \in X$,


$\delta \geq 0 $, $\varepsilon \geq 0$, $N \subset X$,


$\mu =\max (\varepsilon,\delta)$,


$\lambda = \mu + 2xy + 2\alpha (M,N)$. Тогда 


\begin{itemize}


\item[a)] при $\lambda > 0$ выполняется неравенство


$$


\beta (y_{N}^{\delta },x_{M}^{\varepsilon }) \leq \beta (N,M) +


\frac{\beta (x_{M}^{\lambda },x_{M}^{+ 0})}{\lambda }


(|\varepsilon - \delta | + 2xy + 2\alpha (M,N));


$$


\item[b)] при $|\varepsilon - \delta | + 2xy + 2\alpha (M,N) < \mu $


выполняется неравенство


$$


\alpha (y_{N}^{\delta },x_{M}^{\varepsilon }) \leq \alpha (N,M) +


\frac{\beta (x_{M}^{\mu },x_{M}^{+ 0})}{\mu }(|\varepsilon - \delta |


+ 2xy + 2\alpha (M,N)). 


$$


\end{itemize}


\end{lem}





Oбобщает теорему 4 из [2]





\begin{thm}


Пусть $M$ --- выпуклое множество пространства $X$, 


удовлетворяющего условиям {\em A, B, C;} $x,  y \in X$,


$N \subset X$, $\mu =\max (\varepsilon,\delta) > 0$.


Тогда выполняется неравенство


$$


\alpha (y_{N}^{\delta },x_{M}^{\varepsilon }) \leq \alpha (N,M) + 


(\frac{2xM}{\mu } + 2)(|\varepsilon - \delta | + 2xy + 2\alpha (M,N)). 


$$


\end{thm}





Aналогична теореме 5 из [2]





\begin{thm}


Пусть $M$, $N$ --- выпуклые множества пространства $X$, 


удовлетворяющего условиям {\em A, B, C;} $x,  y \in X$,


$\lambda =\max (\varepsilon,\delta) + 2xy + 2\alpha (M,N)$.


Тогда при $\lambda > 0$ выполняется неравенство


\begin{equation}


\alpha (y_{N}^{\delta },x_{M}^{\varepsilon }) \leq \alpha (N,M) + 


\bigg(\frac{2\min (xM,yN)}{\lambda } + 1\bigg)(|\varepsilon - \delta |


+ 2xy + 2\alpha (M,N)).


\end{equation}


\end{thm}





{\bf Замечание.}


Доказательства леммы 4, теорем 1, 2 по существу не отличаются


от соответствующих частей доказательств леммы 10, теорем 4, 5


из [2], если учесть следствие и лемму 3. Эти доказательства


приведены лишь для удобства чтения.





\section{Доказательства полученных результатов}





{\bf Доказательство леммы 2}. a) 


Пусть $\beta (F,x_{M}^{\varepsilon }) > 0$ и число $s > 0$ достаточно мало.


Выберем точку $z \in F \setminus x_{M}^{\varepsilon }$ так, что 


\begin{equation}


zx_{M}^{\varepsilon } \geq \beta (F,x_{M}^{\varepsilon }) - s.


\end{equation}


Затем выберем точку $p \in x_{M}^{t}$ так, что


\begin{equation}


zp \leq zx_{M}^{t} + s.


\end{equation}


Введем обозначения: $a = [x,z]\cap S(x,xM + t)$,


$b = [x,z]\cap S(x,xM + \varepsilon )$, 


$c = [p,z]\cap S(x,xM + \varepsilon )$. В силу (3)


\begin{equation}


zc \geq \beta (F,x_{M}^{\varepsilon }) - s.


\end{equation}


А из (1) следует неравенство


$$


xc \leq px\frac{cz}{pz} + xz\frac{cp}{pz}.


$$


Но $xc = xb = xz - bz$, $ax \geq xp$,


$pc = pz - cz$. Поэтому 


$$


xz - bz \leq ax\frac{cz}{pz} + xz\frac{pz - cz}{pz}.


$$


Отсюда получаем


$cz \leq pz\frac{bz}{az}$.


Кроме того, $\delta - t > \delta - \varepsilon$


и величина $\tau = \delta - \varepsilon - bz =


\delta - t - az$ неотрицательная. Следовательно, 


$$


cz \leq pz\frac{bz}{az} \leq


pz\frac{\delta - \varepsilon - \tau }{\delta - t - \tau } \leq


pz\frac{\delta - \varepsilon}{\delta - t}. 


$$


Из этих неравенств, а также из (4), (5) следует


$$


\beta (F,x_{M}^{\varepsilon }) -s \leq


cz \leq (zx_{M}^{t} + s)\frac{\delta - \varepsilon}{\delta - t} 


\leq (\beta (F,x_{M}^{t}) + s)\frac{\delta - \varepsilon}{\delta - t}.


$$


Переходя к пределу при $s \to 0+$, получим неравенство a).


 


Докажем неравенство b).


Пусть $\beta (x_{M}^{\delta },G) > 0$ и число $s > 0$ 


достаточно мало. Выберем точку 


$z \in x_{M}^{\delta } \setminus G$ так, что 


$zG \geq \beta (x_{M}^{\delta },G) - s$.


Затем выберем такую точку $p \in G$, что


\begin{equation}


zp \leq zG + s.


\end{equation}


Тогда


\begin{equation}


\beta (x_{M}^{\delta },G) -s \leq zp \leq \beta (x_{M}^{\delta },G) + s.


\end{equation}


Если для каждого достаточно малого $s$ точка $z$ 


принадлежит $x_{M}^{\varepsilon}$, то 


$\beta (x_{M}^{\delta },G) = \beta (x_{M}^{\varepsilon },G)$


и неравенство b) верно. Предположим, что 


$z\in x_{M}^{\delta } \setminus x_{M}^{\varepsilon }$ и введем обозначения:


$b=[x,z]\cap S(x,xM + \varepsilon )$,


$c = [p,z]\cap S(x,xM + \varepsilon )$.


В силу (1)


$xc \leq px\frac{cz}{pz} + xz\frac{cp}{pz} =


px\frac{pz - pc}{pz} + xz\frac{cp}{pz} = px + pc\frac{xz - px}{pz}$.


Но $xc = xb$. Следовательно,


\begin{equation}


pz \leq pс\frac{xz - xp}{xb - px} \leq


pc\frac{\delta - t }{\varepsilon - t}.


\end{equation}


Кроме того, из (6) и неравенства треугольника следует


\begin{equation}


cp = pz - cz \leq zG + s - cz \leq cG + s 


\leq \beta (x_{M}^{\varepsilon },G) + s.


\end{equation}


Тогда из (7)--(9) получим


$$


\beta (x_{M}^{\delta },G) -s \leq zp \leq


pc\frac{\delta - t }{\varepsilon - t}


\leq (\beta (x_{M}^{\delta },G) + s)\frac{\delta - t }{\varepsilon - t}.


$$


Переходя к пределу при $s \to 0+$, получим неравенство b).


 


Для доказательства c) достаточно применить неравенство b)


при $G = x_{M}^{\varepsilon '}$, $t = \varepsilon '$,


$\varepsilon = \delta '$,


а затем a) при $F = x_{M}^{\delta }$$.\quad\square$


\medskip





{\bf Доказательство следствия.}


a) Достаточно в неравенстве с) леммы 2 перейти к пределам при


$\varepsilon' \to 0$, $\varepsilon \to 0$. 





b) В том же неравенстве нужно положить $\delta ' = \delta$


и перейти к пределу при $\varepsilon ' \to 0$.


 


c) Следует из b).





d) Следует из b) и простого 


неравенства $\beta (x_{M}^{\delta },x_{M}^{+ 0}) \leq 2xM + \delta$ 


([2], (3.10))$.\quad\square$


\medskip





{\bf Доказательство леммы 4}. a) Из простых оценок ([1], (1.5)) получим 


$$


\delta_{1} = \delta + xy + yN + \beta (N,M) - xM \leq


\delta + 2xy + 2\alpha (M,N) \leq \lambda. 


$$


Тогда из лемм 1, 3 и п.\,b) следствия получим


\begin{multline*}


\beta (y_{N}^{\delta },x_{M}^{\varepsilon }) \leq \beta (N,M) +


\beta (x_{M}^{\lambda },x_{M}^{\varepsilon }) \leq \beta (N,M) +


\frac{\beta(x_{M}^{\lambda },x_{M}^{+0})}{\lambda}(\lambda-\varepsilon)\leq\\


%


\le \beta (N,M) +\frac{\beta (x_{M}^{\lambda },x_{M}^{+ 0})}{\lambda }


(|\varepsilon - \delta | + 2xy + 2\alpha (M,N)).


\end{multline*}





b) Из простых оценок ([1], (1.5)) получим


\begin{multline*}


|\mu - \delta_{2}| = |\mu - \delta - yN + xy + xM + \beta (M,N)| \leq


|\mu - \delta | + 2xy + 2\alpha (M,N) \leq \\


%


\le |\varepsilon - \delta | + 2xy + 2\alpha (M,N) < \mu.


\end{multline*}


Следовательно, $\delta_{2} > 0$ и можно применить


леммы 1, 3 и следствие b). Тогда


\begin{multline*}


\beta (x_{M}^{\varepsilon },y_{N}^{\delta }) \leq \beta (M,N) +


\beta (x_{M}^{\varepsilon },x_{M}^{\delta_{2} - 0}) \leq \beta (M,N) +


\beta (x_{M}^{\mu },x_{M}^{\delta_{2}}) \leq \\


%


\le \beta (M,N) +\frac{\beta (x_{M}^{\mu },x_{M}^{+ 0})}{\mu }


(|\varepsilon - \delta | + 2xy + 2\alpha (M,N)).\quad\square


\end{multline*}





{\bf Доказательство теоремы 1}. Из леммы 4 a) и следствия d) получим


\begin{multline*}


\beta (y_{N}^{\delta },x_{M}^{\varepsilon }) \leq \beta (N,M) +


\frac{\beta (x_{M}^{\lambda },x_{M}^{+ 0})}{\lambda }


(|\varepsilon - \delta | + 2xy + 2\alpha (M,N)) \leq \\


\le\alpha (N,M) + \bigg(\frac{2xM}{\mu } + 1\bigg)(|\varepsilon - \delta |


+ 2xy + 2\alpha (M,N)).


\end{multline*}


А при $|\varepsilon - \delta | + 2xy + 2\alpha (M,N) < \mu $ 


из леммы 4 a) и следствия d) получим


$$


\alpha (y_{N}^{\delta },x_{M}^{\varepsilon }) \leq \alpha (N,M) +


(\frac{2xM}{\mu } + 1)(|\varepsilon - \delta | + 2xy + 2\alpha (M,N)). 


$$





Пусть теперь $|\varepsilon - \delta | + 2xy + 2\alpha (M,N) \geq \mu$. Тогда


\begin{multline*}


\beta (x_{M}^{\varepsilon },y_{N}^{\delta }) \leq


\beta (x_{M}^{\varepsilon },x) + xy + yy_{N}^{\delta } \leq 


xM + \varepsilon + xy + yN \leq \\


\le 2xM + |xM - yN| + \varepsilon + xy \leq


2xM + 2xy + \alpha (M,N) + \varepsilon \leq \\


\le 2xM + \mu + |\varepsilon - \delta | + 2xy + 2\alpha (M,N) \leq


(\frac{2xM}{\mu } + 2)(|\varepsilon - \delta | + 2xy + 2\alpha (M,N)),


\end{multline*}


отсюда следует теорема 1$.\quad\square$


\medskip





{\bf Доказательство теоремы 2}. Обозначим правую часть в (2) 


через $A$. Рассмотрим три случая для оценки выражения 


$\beta (y_{N}^{\delta },x_{M}^{\varepsilon })$.





1. $\delta_{1} = \delta + xy + yN + \beta (N,M) - xM = 0$.





Тогда из неравенства $xM - yN \leq xy + \beta (N,M)$ 


следует, что $\delta = 0$. Из леммы 1 получим


\begin{multline*}


\beta (y_{N}^{\delta },x_{M}^{\varepsilon }) \leq


\beta (N,M) + \beta (x_{M}^{+0},x_{M}^{\varepsilon }) \leq 


\beta (N,M) + 2xM \leq


\beta (N,M) + \\ + 2\min (xM,yN) + 2|xM - yN| \leq


\beta (N,M) + 2\min (xM,yN) + 2xy + 2\alpha (M,N) 


\leq A.


\end{multline*}





2. $0 < \delta_{1} \leq \varepsilon $. 


Тогда из лемм 1, 3 получим


$\beta (y_{N}^{\delta },x_{M}^{\varepsilon }) \leq \beta (N,M) \leq A$.





3. $\delta_{1} > \varepsilon $.


Сначала получим простые неравенства $\delta_{1} < \lambda$, 


\begin{multline*}


2xM+\delta_{1} = \delta + xy + 2\min (xM,yN) + |xM - yN| + \beta (N,M)\leq\\


\le \mu + 2xy


+ 2\min (xM,yN) + 2\alpha (M,N) \leq \lambda + 2\min (xM,yN),


\end{multline*}


$\frac{\delta_{1} - \varepsilon }{\delta_{1}} \leq


\frac{\lambda - \varepsilon }{\lambda }$.


Используя эти неравенства, а также лемму 1 и следствие d), получим


\begin{multline*}


\beta (y_{N}^{\delta },x_{M}^{\varepsilon }) \leq


\beta (N,M) + \beta (x_{M}^{\delta_{1}},x_{M}^{\varepsilon }) \leq 


\beta (N,M) + (2xM + \delta_{1})\frac{\delta_{1} - \varepsilon }


{\delta_{1}} \leq \\ \le \beta (N,M) + 


(2\min (xM,yN) + \lambda) \frac{\lambda - \varepsilon }{\lambda } \leq A.


\end{multline*}


И утверждение теоремы следует из симметричности правой части 


последнего неравенства$.\quad\square$





\begin{thebibliography}{9}





\bibitem{1} Маринов А.В.


{\em Устойчивость $\varepsilon$-квазирешений операторных уравнений


$1$ рода} // Приближение функций полиномами и сплайнами. 


-- Свердловск, 1985. -- C.\,105--117.


\bibitem{2} Маринов А.В.


{\em Оценки устойчивости метрической $\varepsilon$-проекции 


через модуль выпуклости пространства} //


Тр. ин-та матем. и мех. УрО РАН. -- 1992. -- Т.\,2. -- C.\,85--109.


\bibitem{3} Буземан Г.


{\em Геометрия геодезических}. -- М.: Физматгиз, 1962. -- 503 c.


\bibitem{4} Busemann H., Phadke B.B. 


{\em Spaces with distinguished geodesics}. --


New~York--Basel: Marsel Dekker Inc., 1987. -- 159 p.


\bibitem{5} Alexander S.B., Bishop R.L. 


{\em The Hadamard--Cartan theorem in locally convex metric 


spaces} // L'Enseign. Math. -- 1990. -- V.\,36. -- P.\,309--320.





\end{thebibliography}





\vskip 2cm





\post{\it Казанский государственный}{\it Поступила}


\post{\it педагогический университет}{24.12.1999}





\label{end}


\end{document}


